\chapter{Conclusion}
\label{ch:conclusion}

This thesis presented an integrated framework for learning and formally verifying neural imitation controllers using 
learned Lyapunov functions. The proposed approach combines imitation learning based on an expert MPC with 
falsification-guided co-training of a neural control policy and the corresponding Lyapunov candidate, 
followed by formal verification using the $\alpha,\beta$-CROWN framework.

For the linear double integrator, the empirical results demonstrate that imitation learning provides an 
effective inductive prior, enabling compact networks to accurately replicate optimal behavior. 
Building upon this initialization, falsification-guided co-training successfully refines both the controller and 
the neural Lyapunov candidate by systematically reducing mined counterexamples. 
Subsequent formal verification using $\alpha,\beta$-CROWN certifies a region of attraction 
covering $65.07\,\%$ of the considered state domain $\mathcal{X}$ for the closed-loop neural controller. 
While enforcing the Lyapunov decrease condition in this domain causes the policy to deviate slightly
from the expert's behavior, this demonstrates that it is possible to trade imitation accuracy
for formal guarantees of closed-loop stability.
In contrast, the nonlinear cartpole benchmark revealed key limitations of the proposed method. 
Although the co-training procedure converged toward lower condition losses and fewer counterexamples, 
the overall performance of the co-trained controller degraded compared to the imitation baseline. 
Furthermore, formal verification of the co-trained controller failed to yield stability guarantees. 
This highlights that the convergence of the objective alone is insufficient to ensure a verifiable region of attraction.

To address these challenges, future work could focus on targeted regularization terms, for example, to ensure a substantial 
contribution of the feature term to the overall Lyapunov function, effectively preventing feature term collapse. 
Furthermore, the falsification loop could be extended from single-step transitions to multi-step rollout horizons 
to suppress actuator oscillations and penalize cumulative dynamic error growth. 
Additionally, the issue of the curse of dimensionality in the verification pipeline could be addressed by 
changing the verification strategy from input-space partitioning to intermediate layer splitting using $\beta$-CROWN. 
Finally, incorporating bounded disturbance sets and parametric uncertainties directly into the computational 
graph would extend the nominal certificates toward robust sim-to-real transfer.

In summary, this thesis establishes a foundational framework for integrating imitation learning with formal 
neural Lyapunov certification. While providing a promising methodological basis for safety-critical neural control, 
scaling verifiable stability guarantees to complex nonlinear dynamics remains an open challenge that requires 
the co-design of expressive architectures, trajectory-oriented falsification, and verifier-aligned optimization.